\section{结论}
定义延迟信息需求的那个验证器，同时也标注了哪些更早屏能回答它，因此回看可以在任何智能体轨迹语料上训练，而无需任何人做标注。起作用的是这份监督：它把单次调用的全历史读取器抬高 10.6 个点、把逐屏打分器抬高 22.3 个点，并逆转了两种架构孰优孰劣。分解本身再带来 6.9 个点的定位与 2.1 个点的端到端，配对检验对二者都不可区分；而在 8--12 屏的距离上它带来 21.8 个点，那里一次调用必须权衡许多彼此难以区分的候选。在第四批样本上，两个训练变体相对 prompt 都保持 17--20 个点的定位优势（79.6\% 与 76.3\%），端到端达到 49.5\% 与 51.6\%，而 prompt 为 23.7\%，二者之间仍然不可区分。13 屏及以上的距离仍未解决：在那里两个训练变体都低于 prompt 全历史。
